Un tireur à l'arc s'entraîne sur une cible dans le but d'atteindre son centre.

On modélise la situation de la façon suivante :

\begin{itemize}
\item au premier tir, il atteint le centre de la cible avec une probabilité de $\dfrac{1}{2}$ ;
\item pour les tirs suivants :
\begin{itemize}
\item lorsqu'il a atteint le centre de la cible au tir précédent, la probabilité qu'il atteigne à nouveau le centre de la cible est $\dfrac{4}{5}$ ;
\item lorsqu'il n'a pas atteint le centre de la cible au tir précédent, la probabilité qu'il atteigne le centre de la cible est $\dfrac{1}{3}$.
\end{itemize}
\end{itemize}

Pour tout entier naturel $n$ non nul, on considère l'événement $T_n$ : « Le tireur atteint le centre de la cible au $n$-ième tir ».

On note $p_n = P\left(T_n\right)$ la probabilité que l'évènement $T_n$ se réalise.

\begin{enumerate}
\item Donner la valeur de $p_1$ et montrer que $p_2 = \dfrac{17}{30}$.

\item Recopier sur la copie l'arbre de probabilité suivant et compléter les pointillés avec les probabilités qui conviennent :

\begin{center}
\begin{tikzpicture}[x=1cm,y=1cm,>=Latex,thick]
\node (O) at (0,0) {};
\node (Tn) at (3,1.25) {$T_n$};
\node (nTn) at (3,-1.25) {$\overline{T_n}$};
\node (TnT) at (6,2) {$T_{n+1}$};
\node (TnnT) at (6,0.5) {$\overline{T_{n+1}}$};
\node (nTnT) at (6,-0.5) {$T_{n+1}$};
\node (nTnnT) at (6,-2) {$\overline{T_{n+1}}$};
\draw (O) -- node[above left] {$p_n$} (Tn);
\draw (O) -- node[below left] {$\ldots$} (nTn);
\draw (Tn) -- node[above] {$\dfrac{4}{5}$} (TnT);
\draw (Tn) -- node[below] {$\ldots$} (TnnT);
\draw (nTn) -- node[above] {$\ldots$} (nTnT);
\draw (nTn) -- node[below] {$\ldots$} (nTnnT);
\end{tikzpicture}
\end{center}

\item Montrer que pour tout entier naturel $n$ non nul :
\[
p_{n+1}=\dfrac{7}{15}p_n+\dfrac{1}{3}
\]

\item On considère la suite $\left(u_n\right)$ définie pour tout entier naturel $n$ non nul par :
\[
u_n=p_n-\dfrac{5}{8}
\]

\begin{enumerate}
\item Montrer que la suite $\left(u_n\right)$ est une suite géométrique de raison $\dfrac{7}{15}$.
\item Déterminer une expression de $u_n$ en fonction de $n$.
\item En déduire une expression de $p_n$ en fonction de $n$.
\end{enumerate}

\item Déterminer la limite de la suite $\left(p_n\right)$ et interpréter cette limite dans le contexte de l'exercice.

\item On considère ci-dessous une fonction seuil, incomplète, écrite en langage Python.

Recopier cette fonction sur la copie en complétant les pointillés afin qu'elle renvoie la plus petite valeur de l'entier $n$ telle que $p_n$ soit supérieur ou égal à $0,6$.

\begin{Verbatim}[frame=single]
def seuil():
    n = 1
    p = 0.5
    while .........:
        n = .........
        p = .........
    return .........
\end{Verbatim}

\item Résoudre dans $\N$ l'inéquation $p_n > 0,6$.
\end{enumerate}
